\section{Discussion}
The frozen experiment provides no demonstrated improvement from the evaluated fixed-kernel augmentation over worst-query acquisition. That conclusion is local: the post-primary matrix includes a forward rank-four split with a benefit on both budget grids, alongside six of eight point contrasts in the opposite direction. The evidence therefore supports sensitivity and limited reliability of this particular surrogate, rather than a general claim that discrepancy-aware design cannot help. This conditional success is scientifically relevant and remains in the main results despite departing from the original negative pattern.

Several mechanisms explain why a plausible variance term need not improve prediction. Acquisition assumes a weak $.01$ design prior, while the actual predictor can select strong shrinkage or zero. The kernel treats the linear trend as known, has fixed unit amplitude, and is not part of the prediction model. Its posterior variance neither estimates realized nonlinear discrepancy nor optimizes paid-validation ridge risk. Moreover, worst-query targets a maximum whereas the endpoint averages queries. A future comparison could fit discrepancy amplitude using only paid data and align the design criterion with predictive risk, but that would be a new method requiring its own development and frozen evaluation.

The support failure sharpens a separate limitation. Numerical span is necessary for unregularized linear estimability when coordinates are fixed and correct; it is not evidence that a biological intervention response can be transported. Source error can fill a missing direction with noise, and ridge can stabilize an extrapolation without identifying it. Our constructed counterexample and interval failures demonstrate these possibilities but do not measure their real-world frequency. Similarly, full support in the original pilot means its negative result cannot by itself judge acquisition's value during an identification transition.

In the original low-budget replay, a large assessment-noise component, low absolute-direction and ordinal scores, heterogeneous query effects and strong zero/shrinkage references limit practical signal. The new rank-four full-grid reductions versus zero are larger, particularly forward, but squared-error gains alone do not establish useful biological forecasts. The selected 525-target cohort, released gene panels and count-uniform pools restrict the population. Survival, capture, effector and time-point differences remain. Bidirectional and repeated-split replay reuse the same data, and no additional donor or biological replicate is manufactured by resampling. Conditional query, cell and technical-batch ranges answer different questions; none supplies an interval for generalization to new biological contexts. There is no equivalence margin, power argument for absence of benefit, or multiplicity-adjusted search for a winning subgroup.

The resource nevertheless permits precise questions that are often obscured by differing information access. Source data, development exposures, scorer resources and actual recipient purchases are accounted separately. Common backbones isolate allocation from prediction; shared-control covariance links depth and control decisions; all paired contributions and comparator combinations are retained. This auditability supports trust in the stated experiment, but engineering checks cannot substitute for a scientific effect or independent validation. The reconstructed neural comparator and hypothetical action prices further restrict comparisons with published methods and laboratory workflows.

For researchers considering this rule, the evidence does not justify adopting the fixed augmentation as a generally superior default. It motivates testing representation, split, budget and predictor alignment before a prospective campaign, with biological replication and feasible acquisition economics. The reusable result is a narrowly specified evaluation of conditional successes and failures, together with a source-error counterexample and an executable assessment framework.
